\documentclass[12pt]{article}
\usepackage[T1]{fontenc}
\usepackage[utf8]{inputenc}
\usepackage{amsmath,amssymb,mathrsfs}
\usepackage{pst-all,pst-func,pst-3dplot,pst-eucl}
\usepackage{multido,color}
\newcommand{\R}{\mathbb{R}}
\newcommand{\N}{\mathbb{N}}
\newcommand{\D}{\mathbb{D}}
\newcommand{\Z}{\mathbb{Z}}
\newcommand{\Q}{\mathbb{Q}}
\newcommand{\C}{\mathbb{C}}
\newcommand{\vect}[1]{\overrightarrow{\,\mathstrut#1\,}}
\newcommand{\barre}[1]{\overline{\,\mathstrut#1\,}}
\newcommand{\pg}{\geqslant}
\newcommand{\pp}{\leqslant}
\newcommand{\e}{\,\text{e}\,}
\renewcommand{\d}{\,\text{d}}
\newcommand{\Cg}{\texttt{]}}
\newcommand{\Cd}{\texttt{[}}
\def\Oij{$\left(\text{O}\,;\,\vec{\imath},\,\vec{\jmath}\right)$}
\def\Oijk{$\left(\text{O}\,;\,\vec{\imath},\,\vec{\jmath},\,\vec{k}\right)$}
\pagestyle{empty}
\begin{document}
\begin{center}
%\begin{tikzpicture}%compiler avec pgfplots
%\begin{axis}[/pgf/number format/.cd, use comma,
%x={15mm}, y={6mm}, %échelle originale : x=20mm
%xmin=-3, xmax=5, ymin = -8, ymax= 7, %limites du graphique
%xtick = {-2,...,4}, ytick={-7,...,6},
%tick label style={font=\footnotesize},
%minor tick num = 0, grid=major, axis lines =center]
%\node[below left] at (axis cs: 0,0) {{\footnotesize 0}};
%\addplot [line width=1.2pt, color=red, smooth, samples=300, domain= -3:5]{(x^2+3*x+1)*exp(-x)};
%\node[red] at(axis cs: 2.2,2) {$ \mathcal{C}_1 $};
%\addplot [line width=1.2pt, color=blue, dashed, smooth, samples=300, domain= -3:5]{(-x^2-x+2)*exp(-x)};
%\node[blue] at(axis cs: 2.2,-1.1) {$ \mathcal{C}_2 $};
%\end{axis}
%\end{tikzpicture}

%Version en pstricks:
\psset{xunit=1.5cm,yunit=0.6cm} %sur le sujet original : xunit=2cm
\begin{pspicture*}(-3.,-8.)(5.,7.)
% Configuration de la grille et des axes
\psgrid[gridlabels=0pt,subgriddiv=0,gridcolor=lightgray,gridwidth=0.25pt](0,0)(-3,-8)(5,7)
%Axes
\psaxes[linewidth=1.25pt,labelFontSize=\scriptstyle,ticks=xy]{->}(0,0)(-2.99,-7.99)(5,7)
\uput[-135](0,0){\footnotesize 0}
% Courbe C1 : (x^2+3*x+1)*exp(-x)
\psplot[linewidth=1.2pt,linecolor=red,plotpoints=2000]{-3}{5}{x dup mul 3 x mul add 1 add 2.71828 x neg exp mul}
\uput[0](2.2,2){\red $\mathcal{C}_1$}
% Courbe C2 : (-x^2-x+2)*exp(-x)
\psplot[linewidth=1.2pt,linecolor=blue,linestyle=dashed,plotpoints=2000]{-3}{5}{x dup mul neg x sub 2 add 2.71828 x neg exp mul}
\uput[0](2.2,-1.1){\blue $\mathcal{C}_2$}
\end{pspicture*}
\end{center}
\end{document}
